\subsection{Exact checks of two rectification constructions}

We next test concrete ways of converting the opened comparison
theorems into strict PD comparisons.  These statements concern the
specified constructions; the existence questions in Part A are
retained afterwards.

\paragraph{A-S09-06 (Proved; the underived adjoint changes the interval).}
Let $D$ be shuffle denormalization from nonnegative, strictly
graded-commutative cochain algebras over $\mathbb Z$ to cosimplicial
commutative rings.  Put $C=\mathbb Z^{\Delta[1]}$, with pointwise
product.  The functor
\[
 B\longmapsto\operatorname{Hom}_{\mathsf{cRing}}(C,D(B))
\]
is represented in cochain algebras by the degree-zero algebra
$\mathbb Z\times\mathbb Z$.  The analogous functor for the constant
cosimplicial ring $\mathbb Z$ is represented by $\mathbb Z$.
The constant-function quasi-isomorphism
$\mathbb Z\to C$ therefore induces the diagonal
$\mathbb Z\to\mathbb Z\times\mathbb Z$ on these representing objects,
which is not a quasi-isomorphism.  In particular, the underived left
adjoint to shuffle denormalization does not preserve this
quasi-isomorphism.

\begin{proof}
We use the integral formula for shuffle denormalization from
Castiglioni--Corti\~nas, \emph{Cosimplicial versus DG-rings},
Section 4, equation (26), and Product structure 5.3,
equation (49), published pages 124 and 131:
\[
 D^nB=\bigoplus_{i=0}^{n}B^i\otimes\Lambda^i\mathbb Z^n.
 \tag{S09.8}
\]
Its product is the product of $B$ tensored with exterior multiplication
in the displayed matching degrees.  The kernel $J_n$ of the iterated
codegeneracy $D^nB\to B^0$ is the positive-exterior-degree ideal;
$J_n^{n+1}=0$.  Idempotents lift uniquely through a nilpotent ideal
in a commutative ring: two lifts $e,f$ have difference nilpotent,
and the orthogonal decomposition by $e$ reduces $f$ to an idempotent
nilpotent or to one minus such an element, forcing $f=e$.
Existence here is given by the degree-zero summand in (S09.8).

Every idempotent $e\in B^0$ is closed.  Indeed
$de=2e\,de$, and $(1-2e)^2=1$, so $de=0$.  Its images under all
coface composites consequently agree.  Thus idempotents in $D^nB$
are exactly these constant images of the idempotents in $B^0$.

The ring $C^n$ is a finite product of copies of $\mathbb Z$,
one for each nondecreasing map $\alpha:[n]\to[1]$.
A unital ring map $C^n\to D^nB$ sends their characteristic functions
to a complete family of orthogonal idempotents.  For nonconstant
$\alpha$, the characteristic function restricts to zero at both
constant vertices under the iterated codegeneracy.  Its image is
an idempotent in $J_n$, and hence is zero.  The two remaining
idempotents are $e$ and $1-e$, for a unique idempotent $e\in B^0$.
Conversely these assignments define a cosimplicial map: the face or
degeneracy image of a constant simplex is constant, so after killing
all nonconstant characteristic functions every compatibility reduces
to the same assignment on the two vertices.  We have proved a
natural bijection
\[
 \operatorname{Hom}_{\mathsf{cRing}}(C,D(B))
 \cong\operatorname{Idem}(B^0)
 \cong\operatorname{Hom}_{\mathsf{CDGA}_{\geq0}}
       (\mathbb Z\times\mathbb Z,B).
\]
This establishes the representing object without assuming an
adjoint exists on other inputs.

The normalized cochain complex of $C$ is
$\mathbb Z^2\xrightarrow{(a,b)\mapsto b-a}\mathbb Z$,
in degrees zero and one.  It has $H^0=\mathbb Z$ diagonally and
no higher cohomology.  Thus the stated constant-function map is
a quasi-isomorphism.  On representing objects it sends the unit
to $(1,1)$, so its degree-zero cokernel is $\mathbb Z$.
This proves its failure to be a quasi-isomorphism.  The finite
certificate also enumerates idempotents in
$(\mathbb Z/m)[t]/(t^3)$ for $m=2,3,4,6$, checking the nilpotent
idempotent step in 315 elements; the all-ring argument above is
the proof.
\end{proof}

\paragraph{A-S09-07 (Proved; exact products in the opened monoidal constructions).}
The Alexander--Whitney cochain construction and the associative
replacement $Q$ of Castiglioni--Corti\~nas have the following explicit
failures of commutativity.  In normalized integral cochains of
$\Delta[1]$, let $e_0$ be the characteristic function of the initial
vertex and let $u$ take value one on its nondegenerate edge.  Then
\[
 e_0\smile u=u,\qquad u\smile e_0=0.
 \tag{S09.9}
\]
For $B=\Gamma_{\mathbb Z}(x,z)\otimes\Lambda(dx,dz)$, the
two-variable strict PD de Rham algebra, $Q^1B$ contains elements
with commutator
\[
 (dx\otimes v)\circ z-z\circ(dx\otimes v)
 =-dx\,dz\otimes v^2\ne0.
 \tag{S09.10}
\]
Thus these particular intermediate algebras do not satisfy the
strict commutativity required of morphism targets in Part A.

\begin{proof}
The Alexander--Whitney formula evaluates a degree-zero factor at
the first vertex when it is on the left and at the last vertex
when it is on the right.  The values of $e_0$ at those vertices
are respectively one and zero, giving (S09.9).

For (S09.10), Product structure 5.3, equation (48), published
page 131 of Castiglioni--Corti\~nas, gives
\[
 (a\otimes r)\circ(b\otimes s)
 =ab\otimes rs+(-1)^{|r|}a\,db\otimes\theta(r)s,
 \qquad\theta(v)=v^2.
\]
Take $a=dx$, $r=v$, $b=z$, $s=1$ and subtract the reversed
product.  The ordinary terms cancel, $d(dx)=0$, and
$\theta(1)=0$, leaving precisely (S09.10).  It is nonzero because
$dx\,dz$ and $v^2$ are basis vectors in the exterior and tensor
factors.  Giving $x,z,dx,dz$ weight one places the entire witness
in weight two, so no completion enters this computation.  The
certificate implements the product on the 13 integral basis
elements of weights at most two, verifies all 2,197 associativity
triples, and computes this nonzero commutator.  Its coefficient
$-1$ remains nonzero after reduction modulo two or three.
\end{proof}

The opened source of the monoidal acyclic-model comparison is
Guill\'en, Navarro, Pascual and Roig, \emph{Monoidal Functors, Acyclic
Models and Chain Operads}, Theorem 9.2, published page 373;
Section 8.1, pages 371--372, explicitly uses Alexander--Whitney
totalization.  Its input axioms in Section 9.1 require full
simplexwise resolutions, horizontal contractibility in every row,
and free modules over a PID.  We have not supplied a strict PD
replacement of its intermediate product (S09.9).  Castiglioni--Corti\~nas
Theorem 9.6, page 139, provides the associative Quillen equivalence
using $Q$; (S09.10) checks the exact attempted restriction of that
functor to the present strict category.  Neither assertion refutes
those published theorems.

\paragraph{A-S09-08 (Proved; a specified projective model structure is impossible).}
The category of nonnegative augmented strictly graded-commutative
PD cochain DGAs over $\mathbb Z$ has no model structure whose
weak equivalences are the underlying quasi-isomorphisms and whose
fibrations are the degreewise surjections.

\begin{proof}
The initial object is $\mathbb Z$.  Set
\[
 E=\Lambda_{\mathbb Z}(x)\otimes\Gamma_{\mathbb Z}(y),
 \qquad |x|=1,\quad |y|=2,\quad dx=y,\quad dy=0,
\]
with its evident augmentation.  The map $\mathbb Z\to E$ has
the left lifting property against every degreewise surjective
map $B\to C$ of augmented strict PD DGAs.  Indeed a map $E\to C$
is determined by the image $c$ of $x$.  Lift $c$ to $b\in B^1$
using surjectivity and send $x$ to $b$ and $y$ to $db$.
The free PD/exterior property gives an augmented PD-DGA lift.

In any model category a map with the left lifting property against
all fibrations is a trivial cofibration.  For completeness, factor
the map as a trivial cofibration followed by a fibration; its lifting
property makes it a retract of that trivial cofibration, so it is
a weak equivalence.  Our specified classes would therefore make
$\mathbb Z\to E$ a quasi-isomorphism.  But in positive arity $m$
the entire cochain complex of $E$ is
\[
 \mathbb Z x\gamma_{m-1}(y)
 \xrightarrow{\ m\ }\mathbb Z\gamma_m(y)
 \quad\text{in degrees }2m-1,2m.
\]
For $m=2$ this gives $H^4(E)=\mathbb Z/2$, contradicting that
conclusion.  The arity-two matrix is the one-by-one matrix $(2)$;
the attached certificate records this finite integral obstruction.
\end{proof}

This conclusion specifies both proposed model classes.  It does
not exclude constructions with other fibrations, auxiliary algebraic
structures, or explicit replacements that have separately proved
admissibility.  In particular it does not assert anything about a
model structure on the subcategory of admissible theories alone.

